\documentclass[a4paper, 10pt]{article}
\usepackage[catalan]{babel}
\usepackage{tcolorbox}
\usepackage{parskip}
\usepackage{multirow}
\usepackage{graphicx}
\usepackage{caption}
\captionsetup{labelfont=bf}
\usepackage{hyperref}
\usepackage[arrowdel]{physics}
\newtcolorbox{definició}[1]{title=#1}
\usepackage[utf8]{inputenc}
\usepackage[left=1.95cm, right=1.95cm, top=30mm, bottom=20mm]{geometry} 
\usepackage{fancyhdr}
\usepackage{amssymb}
\usepackage{enumitem}
\usepackage{graphicx}
\usepackage{subfigure}
\usepackage{cancel}
\usepackage{float}
\usepackage{titling} % Per controlar el títol
\pagestyle{fancy}
\fancyhf{}
\lhead[\leftmark]{{Universitat Autònoma de Barcelona}}
\rhead{Termodinàmica i Mecànica Estadística}
\cfoot[]{\thepage}
\renewcommand{\labelenumi}{\textbf{\theenumi.}} % Numeració en Bf.

% Redueixo l'espai abans del títol
\setlength{\droptitle}{-2cm}

\begin{document}
	\title{\Huge{\textbf{Segon Lliurament}} \\ \vspace{0.2cm} \large{\textbf{Termodinàmica i Mecànica Estadística}}}
	
	\author{\small{1549086: Bujones Umbert, Jun Shan} \\
		\small{1672980: González Barea, Eric} \\
		\small{1644841: Vilarrúbias Morral, Natàlia}}
	\date{Desembre 2025}
	\maketitle
	
	% Aseguramos que las cabeceras se muestren también en la primera página
	\thispagestyle{fancy}
	
Un objecte de massa $m$ es troba subjecte entre dues molles de constant $k$ i longitud natural
$l$. La distància entre les parets és $2l$ (com es mostra a la figura \ref{fig1}) i la partícula només es pot moure
unidimensionalment. El Hamiltonià de la partícula és
\begin{equation}
	H = \frac{p^2}{2m}+kx^2. \label{eq:enunciat}
\end{equation}

\begin{figure}[h]
	\centering
	\includegraphics[width=0.5\linewidth]{screenshot001}
	\caption{Sistema d'estudi.}
	\label{fig1}
\end{figure}

\section*{Apartat a)}
\textbf{Calculeu el nombre de microstats de la partícula en funció de la seva energia $E$.}

\section*{Apartat b)}
\textbf{Si tot el sistema es troba dins d'un bany tèrmic a temperatura $T$ la partícula estarà sotmesa a un soroll tèrmic. Calculeu la funció de partició de la partícula i el valor mig de l'energia de la partícula $\langle E \rangle$. Establiu els límits d'altes i baixes temperatures i trobeu $\langle E \rangle$ en aquests límits. Raoneu físicament els resultats. Calculeu el valor mig de la posició de la partícula i la variança de la posició i analitzeu i raoneu els límits d'altes i baixes temperatures.}

\begin{enumerate}
	\item \textbf{\textit{Càlcul de la funció de partició.}}

La funció de partició que hem de calcular és la funció de partició canònica $Z$ (tenim el sistema en un bany tèrmic a temperatura $T$). Per calcular-la podem usar la seva definició (que vam veure a teoria) aplicada a aquest cas en què tenim una única partícula sotmesa al hamiltonià de \eqref{eq:enunciat}. Així doncs, tenim que

\begin{equation}
	Z = \int \frac{\mathrm{d}x\mathrm{d}p}{h}e^{-\beta H(x,p)} = \int \frac{\mathrm{d}x\mathrm{d}p}{h}e^{-\beta p^2/2m}e^{-\beta k x^2} = \frac{1}{h} \int_{-l}^{l} e^{-\beta kx^2}\mathrm{d}x\int_{-\infty}^{+\infty} e^{-\beta p^2/2m} \mathrm{d}p. \label{eqb1:1}
\end{equation}
	
Calculem ara per separat les dues integrals anteriors. La segona és una integral Gaussiana estàndard que es pot calcular si recordem que
\begin{equation}
	\int_{-\infty}^{+\infty} e^{-ax^2}\mathrm{d}x = \sqrt{\frac{\pi}{a}}, \label{eqb1:2}
\end{equation}
de manera que tenim que
\begin{equation}
	\int_{-\infty}^{+\infty} e^{-\beta p^2/2m} \mathrm{d}p = \sqrt{\frac{2\pi m}{\beta}}. \label{eqb1:3}
\end{equation}
Per altra banda, la primera integral és també una Gaussiana però amb els límits d'integració diferents de $\pm \infty$ i, per tant, la podem deixar en termes de la \textit{error function} erf$(x)$. Recordem que aquesta funció és la següent
\begin{equation}
	\mathrm{erf}(x) = \frac{2}{\sqrt{\pi}} \int_0^x e^{-t^2}\mathrm{d}t. \label{eqb1:4}
\end{equation}
Si desenvolupem la primera integral de \eqref{eqb1:1} tenim el següent
\begin{equation*}
	\int_{-l}^{l} e^{-\beta kx^2}\mathrm{d}x = 2 \int_{0}^{l} e^{-\beta kx^2}\mathrm{d}x = \frac{2}{\sqrt{\beta k}}\int_0^{l\sqrt{\beta k}}e^{-t^2}\mathrm{d}t = \frac{2}{\sqrt{\beta k}} \cdot \frac{\sqrt{\pi}}{2}\mathrm{erf}(l\sqrt{\beta k}) \Rightarrow
\end{equation*}
\begin{equation}
	\Rightarrow \sqrt{\frac{\pi}{\beta k}} \mathrm{erf}(l\sqrt{\beta k}), \label{eqb1:5}
\end{equation}
on al primer pas hem fet usat que la $e^{-\beta k x^2}$ és una funció simètrica respecte del 0 i, al segon pas, hem fet el següent canvi de variable
\begin{align*}
	\sqrt{\beta k} x &= t, \\
	\sqrt{\beta k} \mathrm{d}x &= \mathrm{d}t.
\end{align*}
Així doncs, si usem els resultats \eqref{eqb1:3} i \eqref{eqb1:5} sobre \eqref{eqb1:1}, trobem la funció de partició canònica pel sistema que estem estudiant
\begin{equation}
	\boxed{Z(\beta) = \frac{\pi}{h\beta}\sqrt{\frac{2m}{k}}\mathrm{erf}(l\sqrt{\beta k})}. \label{eqb1:6}
\end{equation}

Una forma alternativa de calcular aquesta funció de partició és usant que 
\begin{equation}
	Z = \int_{0}^{\infty}\mathrm{d}Ee^{-\beta E}\Omega(E,V,N),
\end{equation}
aprofitant que, a l'apartat anterior, hem calculat el nombre de microestats $\Omega$.
\\
	\item \textbf{\textit{Càlcul de $\langle E \rangle$ i límits d'altes i baixes temperatures.}}
Per calcular $\langle E \rangle$ (que es correspondria amb l'energia interna $U$ un cop prenguéssim el límit termodinàmic $N \rightarrow \infty$), usem que, tal i com es va demostrar a les classes teòriques
\begin{equation}
	\langle E \rangle = - \frac{\partial \ln Z}{\partial \beta}. \label{eqb2:1}
\end{equation}
Si prenem el logaritme neperià de $Z(\beta)$ (equació \eqref{eqb1:6}), tenim que
\begin{equation}
	\ln Z = \ln\left(\frac{1}{\beta}\right) + \ln\left(\frac{\pi}{h}\cdot \sqrt{\frac{2m}{k}}\right) + \ln(\mathrm{erf}(l\sqrt{\beta k})),\label{eqb2:2}
\end{equation}
de manera que
\begin{equation*}
	\frac{\partial \ln Z}{\partial \beta} = \frac{1}{1/\beta}\cdot \left(-\frac{1}{\beta^2}\right) + \frac{\mathrm{d}}{\mathrm{d} \beta}\left[\ln\left(\mathrm{erf}(l\sqrt{\beta k})\right)\right] = -\frac{1}{\beta} + \frac{1}{\mathrm{erf}(l\sqrt{\beta k})}\cdot \frac{\mathrm{d}}{\mathrm{d} \beta}\left[\mathrm{erf}(l\sqrt{\beta k})\right]. 
\end{equation*}
Per fer la darrera derivada, apliquem la regla de la cadena
\begin{equation*}
	\frac{\mathrm{d}}{\mathrm{d}\beta}\left[\mathrm{erf}(l\sqrt{\beta k})\right] = \frac{\mathrm{d}(\mathrm{erf}(l\sqrt{\beta k}))}{\mathrm{d}(l\sqrt{\beta k})}\cdot \frac{\mathrm{d}(l\sqrt{\beta k})}{\mathrm{d}\beta} = \frac{l\sqrt{k}\beta^{-1/2}}{2}\cdot \frac{\mathrm{d}(\mathrm{erf}(l\sqrt{\beta k}))}{\mathrm{d}(l\sqrt{\beta k})} = \ldots
\end{equation*}
usant la definició de la \textit{error function} donada a \eqref{eqb1:4}, obtenim
\begin{equation*}
	\ldots = l \sqrt{\frac{k}{\pi \beta}} \frac{\mathrm{d}}{\mathrm{d}(l\sqrt{\beta k})}\left[\int_0^{l\sqrt{\beta k}}e^{-t^2}\mathrm{d}t\right] = \ldots
\end{equation*}
i ara, recordant la \textit{regla de Barrow}, arribem a que
\begin{equation*}
	\ldots = l \sqrt{\frac{k}{\pi \beta}}\frac{\mathrm{d}}{\mathrm{d}(l\sqrt{\beta k})}\left[F(l\sqrt{\beta k})- F(0)\right] = l \sqrt{\frac{k}{\pi \beta}}\frac{\mathrm{d}}{\mathrm{d}(l\sqrt{\beta k})}\left[F(l\sqrt{\beta k})\right] = l \sqrt{\frac{k}{\pi \beta}} e^{-l^2\beta k}.
\end{equation*}
Així doncs, si agrupem tot això
\begin{equation}
	\frac{\partial \ln Z}{\partial \beta} = -\frac{1}{\beta} + \sqrt{\frac{kl^2}{\pi \beta}} \cdot \frac{e^{-l^2k\beta}}{\mathrm{erf}(l\sqrt{\beta k})}. \label{eqb2:3}
\end{equation}
Per tant, el valor esperat de l'energia de la partícula serà
\begin{equation}
	\boxed{\langle E \rangle = \frac{1}{\beta} - \sqrt{\frac{kl^2}{\pi \beta}} \cdot \frac{e^{-l^2k\beta}}{\mathrm{erf}(l\sqrt{\beta k})}}. \label{eqb2:4}
\end{equation}

Estudiem ara els límits d'altes i baixes temperatures. Per fer això, hem de comparar l'energia tèrmica del sistema $k_BT \equiv \frac{1}{\beta}$ amb l'energia potencial elàstica $kl^2$ deguda a la presència de les molles. És a dir, hem d'estudiar el comportament del factor
\begin{equation}
	x \equiv \frac{kl^2}{k_BT} = l^2\beta k.
\end{equation}
\begin{enumerate}
	\item \textbf{Comportament de $\langle E \rangle$ a altes temperatures.} En aquest cas tenim que $T \rightarrow \infty$ i, aleshores
	\begin{equation*}
		k_BT \gg kl^2 \Rightarrow 1 \gg \frac{kl^2}{k_BT} \Rightarrow 1 \gg l^2\beta k
	\end{equation*}
	Per veure el comportament de la \textit{error function} quan la variable $l^2k\beta \rightarrow$ usem la seva expansió en sèrie de Taylor (treballant fins a primer ordre)\footnote{Aquest resultat es pot trobar desenvolupant $e^{-z^2}$ en sèrie de Taylor i integrant terme a terme. Veure \url{https://en.wikipedia.org/wiki/Error_function}.}
	\begin{equation}
		\mathrm{erf}(z) = \frac{2}{\sqrt{\pi}}\left(z-\frac{z^3}{3}+\frac{z^5}{10}-\frac{z^7}{42}+\ldots\right) \approx \frac{2}{\sqrt{\pi}}z. \label{eqb2:5}
	\end{equation}
	Això, en el nostre cas, es tradueix en què
	\begin{equation}
		\mathrm{erf}(l\sqrt{\beta k}) \approx 2\sqrt{\frac{l^2\beta k}{\pi}}. \label{eqb2:6}
	\end{equation}
	Per altra banda, per $1 \gg l^2\beta k$, tenim que
	\begin{equation}
		e^{-l^2\beta k} \approx 1. \label{eqb2:7}
	\end{equation}
	Si substituïm les aproximacions per altes temperatures \eqref{eqb2:6} i \eqref{eqb2:7} sobre \eqref{eqb2:5}, trobem que
	\begin{equation*}
		\langle E \rangle = \frac{1}{\beta} - \frac{1}{2\beta} =\frac{1}{\beta} \Rightarrow
	\end{equation*}
	\begin{equation}
		\boxed{\langle E \rangle_{T\rightarrow \infty} \approx \frac{1}{2} k_BT}. \label{eqb2:8}
	\end{equation}
	
	\textbf{RAONAR FÍSICAMENT A PARTIR DEL TEOREMA D'EQUIPARTICIÓ}
	
	\item \textbf{Comportament de $\langle E \rangle$ a baixes temperatures.}
	
	\textbf{AAAAAAAAAAA}
\end{enumerate}



	\item \textbf{\textit{Càlcul de $\langle x \rangle$, $\langle x^2 \rangle$ i $(\Delta x)^2$ i límits d'altes i baixes temperatures.}}
\end{enumerate}

\end{document}